\section{Mathematical and Implementation Details}
\label{app:math}
\setcounter{equation}{0}
\setcounter{table}{0}
\setcounter{figure}{0}

\subsection{Role-Specific Reliability and Training-Only Anchors}

For sample $i$, modality $m$, and role $k\in\{\mathrm{str},\mathrm{num}\}$, the quality/availability indicator is $a_{i,m}=\mathbf{1}[\mathrm{mask}_{i,m}q_{i,m}>0]$. The learned reliability $r_{i,m,k}$ is mapped to the unnormalized routing weight
\begin{equation}
\widetilde w_{i,m,k}=a_{i,m}r_{i,m,k}.
\label{eq:appendix-reliability-weight}
\end{equation}
Routing weights are normalized across available modalities with the denominator clamped at $10^{-6}$ and remain zero if no modality is available. Inference does not require target-derived reliability signals.

The evaluated reliability-prior term is
\begin{equation}
\mathcal L_{\mathrm{rel}}
=
\operatorname{mean}_{a_{i,m}=1}
\left[
(r_{i,m,k}-p_{m,k})^2
\right],
\end{equation}
with fixed $(\mathrm{structural},\mathrm{numerical})$ priors
\begin{table}[t]
\centering
\caption{Fixed role-dependent reliability priors for the evaluated modalities.}\label{tab:appendix-reliability-priors}
\begin{tabular}{lcc}
\toprule
Modality & Structural & Numerical \\
\midrule
Migrated image & 0.85 & 0.45 \\
Horizon & 0.95 & 0.10 \\
RMS velocity & 0.55 & 0.90 \\
Well & 0.25 & 0.95 \\
\bottomrule
\end{tabular}
\end{table}

Training-only anchors are constructed from a one-level Haar decomposition of the target velocity. The numerical anchor is the resized LL component. The structural anchor is the $\ell_2$ fusion of resized LH, HL, and HH components. Haar boundary handling is replicate/edge, resizing is nearest-neighbor, and anchors are detached before condition-loss use.

The subset-Symile terms group samples by identical modality-availability pattern. Within each group, same-record aligned projected modality representations together with the corresponding projected anchor form the positive tuple; negatives are generated by within-group batch derangements, using up to $\mathrm{batch\ size}-1$ negatives. Groups containing fewer than two samples are skipped. Anchor calibration uses $\ell_2$-normalized projected modality representations and normalized projected anchors in a symmetric InfoNCE objective; a single-item group falls back to cosine distance.

Two separation terms are active. Structural/numerical orthogonality is the absolute cosine similarity between adaptively pooled structural and numerical role maps. Unique-role separation averages the absolute cosine similarities between the pooled unique representation and the two pooled role representations. No additional within-role auxiliary loss is active.

\subsection{Gradient Paths and Active Objective Boundaries}

In the prediction-relative residual branch, $\bar B=\operatorname{sg}[\hat B]$ retains the forward background value while preventing residual-loss gradients from updating the background estimator through $\hat B$. The background estimator receives its direct supervision from $\mathcal L_{\mathrm{bg}}$, while shared condition-encoder parameters are jointly optimized through the active objective in Eq.~\eqref{eq:condition-objective}. The low-pass operator $P_L$ is the $5\times5$ stride-one padded average pool described in Section~\ref{sec:methods-background}, with $P_H=I-P_L$.

Pairwise structural/numerical penalties, frequency penalties, a residual quality/rho penalty, and background high-frequency/multiscale/quality-gate penalties have zero active weight in the evaluated configuration. No additional within-role auxiliary separation term is active; modality dropout is disabled.
